\documentclass[journal]{IEEEtran}
\usepackage{cite}
\usepackage{amsmath,amssymb,amsthm}
\usepackage{graphicx}
\usepackage{bm}
\usepackage{epstopdf}

\newcommand{\getname}[1]{eps/#1.eps}

\newcommand{\real}{\mathbb{R}}
\newcommand{\sinc}{\mathrm{sinc}\,}
\newcommand{\bfomega}{\boldsymbol{\omega}}
\newcommand{\bfOmega}{\boldsymbol{\Omega}}
\newcommand{\adj}[1]{\mathrm{ad}({#1})}
\newcommand{\sfSigma}[1]{\mathsf{\Sigma}_\mathnormal{#1}}
\newcommand{\Jg}{\mathsf{J}_\mathnormal{g_r}}
\newcommand{\Jom}{\mathsf{J}_\mathnormal{\omega_r}}

\newcommand{\beq}{\begin{equation}}
\newcommand{\eeq}{\end{equation}}
\newcommand{\beqar}{\begin{eqnarray}}
\newcommand{\eeqar}{\end{eqnarray}}

\newtheorem{theorem}{Theorem}
\newtheorem{proposition}{Proposition}
\newtheorem{corollary}{Corollary}
\newtheorem{remark}{Remark}

\begin{document}
\title{Higher-Order Uncertainty Propagation and Saddlepoint Marginalization on the SE(3) Lie Group}

\author{\IEEEauthorblockN{Frank~O.~Kuehnel}
    \IEEEauthorblockA{Excel Solutions LLC\\
        Email: kuehnel@excelsolutionsllc.com}
    \thanks{Manuscript received XXXX; revised XXXX.}}

\markboth{IEEE Transactions on XXXX, Vol.~XX, No.~X, MONTH~YEAR}%
{Kuehnel: Higher-Order Uncertainty Propagation on SE(3)}

\maketitle

\begin{abstract}
    Bayesian inference over the Special Euclidean group $SE(3)$ requires propagating
    uncertainty through the non-commutative group composition and marginalizing over
    latent variables that enter non-linearly through the projective observation model.
    We develop a self-contained framework addressing both problems.
    First, we establish that Lie-Cartan exponential coordinates of the first kind
    yield unbiased estimates under Gaussian measurement noise, and we provide a
    geometric explanation rooted in the semi-direct product structure of $SE(3)$.
    Second, we derive closed-form composition Jacobians---including the
    rotation-translation coupling Jacobian---that enable exact first-order
    and systematic second-order uncertainty propagation on $SE(3)$,
    going beyond the standard linearized approximation used in current
    Kalman-type filters.
    Third, we introduce a saddlepoint marginalization scheme for eliminating
    latent variables (such as 3D landmark positions) from the joint posterior
    when the observation model is projective and hence non-Gaussian in the
    pose parameters. The saddlepoint approximation yields a
    marginal posterior over $SE(3)$ that is more accurate than the
    Laplace approximation commonly used in bundle adjustment.
    We provide complete derivations of all required Jacobians over $SE(3)$,
    including a novel closed-form formula for composing finite Rodrigues
    vectors via the $SU(2)$ double cover with careful treatment of
    phase reflection at $\Theta = \pi$. Numerical experiments
    demonstrate the bias reduction, improved convergence basins,
    and accuracy gains of the proposed methods.
\end{abstract}

\begin{IEEEkeywords}
    Lie groups, SE(3), uncertainty propagation, saddlepoint approximation,
    Bayesian inference, pose estimation, Rodrigues vector, Baker-Campbell-Hausdorff
\end{IEEEkeywords}

%% ====================================================================
\section{Introduction}
\label{sec:intro}
%% ====================================================================

Estimating the pose of a rigid body from noisy measurements is a
fundamental problem across robotics, computer vision, aerospace
navigation, and surgical guidance. The six-dimensional pose---three
for position, three for orientation---lives on the Special Euclidean
group $SE(3)$, a non-commutative Lie group whose curved geometry
complicates standard probabilistic inference.

\subsection{Motivation and Problem Statement}

Current state-of-the-art methods for pose estimation---whether
filtering-based (Extended Kalman Filters, Unscented Kalman Filters)
or optimization-based (bundle adjustment, factor graph SLAM)---rely
on two simplifying assumptions that limit their accuracy:

\begin{enumerate}
    \item \emph{First-order linearization of uncertainty propagation.}
          When composing uncertain poses $\hat{f} = f_0 \cdot \Delta f$
          where $\Delta f$ is a random perturbation in exponential coordinates,
          the covariance is transported using a first-order Jacobian.
          On $SE(3)$, the non-trivial coupling between rotational and
          translational degrees of freedom (the semi-direct product structure)
          introduces systematic bias at first order that is often neglected.

    \item \emph{Gaussian approximation of inherently non-Gaussian posteriors.}
          When observations are projective (camera measurements),
          the joint posterior over pose and 3D landmarks is non-Gaussian
          in the pose parameters. Standard Laplace approximations
          (equivalent to Gauss-Newton in optimization) capture only the
          mode and curvature, missing asymmetry and heavier tails.
\end{enumerate}

This paper addresses both limitations within a unified $SE(3)$
framework, providing:

\begin{itemize}
    \item Exact first-order and systematic second-order mean and
          covariance propagation formulas for $SE(3)$ composition,
          derived from a closed-form Baker-Campbell-Hausdorff (BCH)
          formula via the $SU(2)/\mathbb{Z}_2 \cong SO(3)$ double cover.

    \item A saddlepoint approximation for marginalizing latent
          Euclidean variables (3D landmarks) from the joint posterior,
          yielding a more accurate marginal over $SE(3)$ than the
          standard Laplace/Gauss-Newton approach.

    \item Complete, closed-form expressions for all required
          $SE(3)$ Jacobians, including the rotation-translation
          coupling Jacobian $\mathsf{J}_t(\bfOmega, \mathbf{t})$,
          with careful treatment of the phase reflection singularity
          at $\Theta = \pi$.
\end{itemize}

\subsection{Related Work}

The use of Lie group structure for pose estimation has a long
history. Pennec and Thirion~\cite{pennec97} introduced
compositional noise models on $SE(3)$. The invariant EKF
(IEKF)~\cite{barrau2017invariant} showed that group-affine
observation models yield autonomous error dynamics.
Sol\`{a} et al.~\cite{sola2018micro} provided a widely-used
reference for Lie group operations in robotics.
Barfoot~\cite{barfoot2024} gave a comprehensive textbook treatment.
Forster et al.~\cite{forster2017manifold} developed
IMU preintegration on manifolds using first-order propagation.

On uncertainty propagation specifically,
Chirikjian~\cite{chirikjian2012stochastic} developed Fourier
analysis on Lie groups for global distributions.
Ye and Chirikjian~\cite{ye2024uncertainty} recently addressed
propagation via stochastic differential equations on unimodular
groups. Wang and Chirikjian~\cite{wang2008nonparametric}
explored non-parametric approaches.
Bourmaud et al.~\cite{bourmaud2013discrete} proposed a
discrete extended Kalman filter on Lie groups.

The present work extends the author's earlier
contributions~\cite{kuehnel2004,kuehnel2005,kuehnel2008}.
In~\cite{kuehnel2004}, the role of Lie-Cartan exponential
coordinates for unbiased pose estimation and the
marginalized MAP estimator were first introduced in
the context of Bayesian inference on $SE(3)$.
In~\cite{kuehnel2005}, the graphical model structure
of the SLAM problem was addressed via local frame
junction trees exploiting the frame dependence of
observations. The present paper refocuses and extends
that foundational work: the second-order propagation
formulas and saddlepoint marginalization are new.

Despite this body of work, several gaps remain:
(i)~closed-form second-order propagation formulas for the
full $SE(3)$ group (not just $SO(3)$) are not available
in the literature;
(ii)~the saddlepoint approximation for marginalizing
Euclidean latent variables from an $SE(3)$ posterior has
not been developed;
(iii)~the detailed algebra connecting the $SU(2)$ double
cover to finite Rodrigues vector composition with phase
reflection handling is scattered and incomplete.

\subsection{Paper Organization}

Section~\ref{sec:SE3geometry} establishes the $SE(3)$
framework and exponential coordinates.
Section~\ref{sec:uncertainty} develops the parametric
uncertainty representation and derives the propagation
formulas to second order.
Section~\ref{sec:saddlepoint} presents the saddlepoint
marginalization scheme.
Section~\ref{sec:experiments} provides numerical validation.
Section~\ref{sec:conclusion} concludes.
The Appendices contain complete derivations of the $SE(3)$
algebra, BCH formula, composition Jacobians, and the
$SU(2)/\mathbb{Z}_2$ connection.

%% ====================================================================
\section{Geometry of $SE(3)$ and Exponential Coordinates}
\label{sec:SE3geometry}
%% ====================================================================

\subsection{The $SE(3)$ Group and Its Lie Algebra}

A rigid body pose $f \in SE(3)$ is represented as
\beq
f \cong (\mathsf{R}(\bfOmega_f), \mathbf{T}_f),\quad
\mathsf{R} \in SO(3),\; \mathbf{T} \in \real^3,
\label{eq:SE3element}
\eeq
with group composition (semi-direct product)
\beq
h \cdot f = (\mathsf{R}_h \mathsf{R}_f,\; \mathsf{R}_h \mathbf{T}_f + \mathbf{T}_h).
\label{eq:groupcomp}
\eeq
The action on 3D points is $f \star \mathbf{x} = \mathsf{R}_f \mathbf{x} + \mathbf{T}_f$.

The Lie algebra $\mathfrak{se}(3)$ has generators
$\{\mathsf{H}_i, \mathsf{T}_j\}$ satisfying
\beq
[\mathsf{H}_i, \mathsf{H}_j]_c = \varepsilon_{ijk}\mathsf{H}_k,\quad
[\mathsf{H}_i, \mathsf{T}_j]_c = \varepsilon_{ijk}\mathsf{T}_k,\quad
[\mathsf{T}_i, \mathsf{T}_j]_c = 0.
\label{eq:sealgebra}
\eeq
The translations form an abelian ideal, making $SE(3)$ a
semi-direct product $SO(3) \ltimes \real^3$.

\subsection{Exponential Map and Lie-Cartan Coordinates}

The exponential map
\beq
\exp: \real^6 \to SE(3),\quad
[\bfOmega, \mathbf{t}] \mapsto
\exp(\hat{\mathbf{H}} \cdot \bfOmega + \hat{\mathbf{T}} \cdot \mathbf{t})
\label{eq:expmap}
\eeq
defines \emph{Lie-Cartan coordinates of the first kind}
$[\bfOmega, \mathbf{t}] \in \real^6$.
The explicit formulas are (Rodrigues):
\beqar
\mathsf{R}(\bfOmega) &=& \mathsf{I} + \sin\Theta\, \mathsf{H}
+ (1-\cos\Theta)\mathsf{H}^2, \label{eq:rodrigues} \\
\mathbf{T} &=& \mathsf{S}(\bfOmega)\, \mathbf{t},\quad
\mathsf{S} = \mathsf{I} + \mathsf{H}^2 +
\frac{1}{\Theta}(\mathsf{I} - \mathsf{R})\mathsf{H}, \nonumber
\eeqar
with $\Theta = |\bfOmega|$, $\mathbf{r} = \bfOmega/\Theta$,
$\mathsf{H} = \hat{\mathbf{H}} \cdot \mathbf{r}$.

\subsection{Why Exponential Coordinates: Bias Analysis}
\label{subsec:bias}

A key observation motivating this work is that the choice of
parametrization affects estimator bias. Consider a camera
localization problem: given matched 3D point sets
$\{\mathbf{x}^j\}$ and $\{\mathbf{p}^j\}$, we solve
\beq
f_0 = \arg\min_{f \in SE(3)} \sum_j
|\mathsf{R}(\bfOmega_f)\mathbf{x}^j + \mathbf{T}_f - \mathbf{p}^j|^2.
\label{eq:localization}
\eeq
When the landmarks are corrupted by Gaussian noise
$\mathbf{x}^j \sim \mathcal{N}(\bar{\mathbf{x}}^j, \sigma)$,
sampling the estimator distribution reveals:

\begin{proposition}[Bias in second-kind coordinates]
    \label{prop:bias}
    If the estimator distribution is unbiased in exponential
    coordinates, $\langle [\bfOmega, \mathbf{t}] \rangle = 0$,
    with non-vanishing rotation-translation correlations
    $\langle \Omega_\alpha t_\beta \rangle \neq 0$, then
    the translational vector $\mathbf{T} = \mathsf{S}(\bfOmega)\mathbf{t}$
    is biased:
    \beq
    \langle \mathbf{T} \rangle =
    \langle \mathsf{S}(\bfOmega)\mathbf{t} \rangle \neq 0.
    \label{eq:biasmechanism}
    \eeq
\end{proposition}

This bias arises because $\mathsf{S}(\bfOmega)$ couples
the rotational and translational components non-linearly.
In second-kind coordinates $[\bfOmega, \mathbf{T}]$,
this coupling is hidden, producing apparent bias.
In first-kind (exponential) coordinates $[\bfOmega, \mathbf{t}]$,
the coupling is explicit and the distribution is centered.

%% ====================================================================
\section{Parametric Uncertainty on $SE(3)$}
\label{sec:uncertainty}
%% ====================================================================

\subsection{Gaussian Model in Exponential Coordinates}

We define the probability of finding pose $f = f_0 \cdot \Delta f$
via the probability of the right deviation
$\ln \Delta f \in \mathfrak{se}(3) \cong \real^6$:
\beq
p_{f_0}(f) \equiv p(\ln f_0^{-1} \cdot f) \sim
\mathcal{N}(\mathbf{0}, \sfSigma{ff_r}),
\label{eq:gaussianSE3}
\eeq
\beq
p(\ln \Delta f) \propto
\exp\!\left(-\frac{1}{2}[\bfOmega, \mathbf{t}]\,
\sfSigma{ff_r}^{-1}\, [\bfOmega, \mathbf{t}]^T\right).
\label{eq:gaussexplicit}
\eeq
We write the uncertainty pose compactly as
$\hat{f} = (f_0, \sfSigma{ff_r})$.

The left-right distinction matters: composing
$\Delta f$ on the left gives a related distribution
\beq
f = \Delta f \cdot f_0,\quad
\sfSigma{ff_l}^{-1} = \mathsf{J}_{lr}^T(f_0)\,
\sfSigma{ff_r}^{-1}\, \mathsf{J}_{lr}(f_0),
\label{eq:leftright}
\eeq
where $\mathsf{J}_{lr}(f_0) \in \real^{6 \times 6}$ is the
adjoint Jacobian (see Appendix~\ref{app:weinorman}).

\subsection{First-Order Composition of Uncertainty Poses}

The composition of two uncertainty poses
$\hat{g} = (g_0, \sfSigma{gg_r})$ and
$\hat{f} = (f_0, \sfSigma{ff_r})$ yields,
at first order:
\beq
\begin{split}
    \hat{h} = \hat{g} \circ \hat{f} \;\longrightarrow\;
    h_0            & = g_0 \cdot f_0,   \\
    \sfSigma{hh_r} & = \sfSigma{ff_r} +
    \mathsf{J}_{lr}(f_0)\, \sfSigma{gg_r}\,
    \mathsf{J}_{lr}^T(f_0).
\end{split}
\label{eq:firstordercomp}
\eeq

\begin{remark}
    In contrast to $\real^n$ where covariances add,
    on $SE(3)$ the ``transported'' covariance
    $\mathsf{J}_{lr}(f_0)\sfSigma{gg_r}\mathsf{J}_{lr}^T(f_0)$
    depends on the base point $f_0$---the Mahalanobis metric
    is position-dependent on the manifold.
\end{remark}

\subsection{Second-Order Corrections}
\label{subsec:secondorder}

The first-order formula~\eqref{eq:firstordercomp} neglects
the curvature of $SE(3)$. We now derive corrections.

Consider the composition $h = g \cdot f$ where both
$g$ and $f$ carry uncertainty. In exponential coordinates,
the composed element has coordinates given by the
BCH formula (Appendix~\ref{app:BCH}):
\beq
\ln(h) = \ln(g \cdot f) =
\bfOmega_g + \bfOmega_f +
\frac{1}{2}[\bfOmega_g, \bfOmega_f]_c + \cdots
\eeq
The second-order contribution to the \emph{mean} is:
\beq
\langle \ln h \rangle =
\langle \bfOmega_g \rangle + \langle \bfOmega_f \rangle +
\frac{1}{2} \langle [\bfOmega_g, \bfOmega_f]_c \rangle +
\mathcal{O}(\sigma^3).
\label{eq:secondordermean}
\eeq

\begin{theorem}[Second-order mean shift]
    \label{thm:meanshift}
    For zero-mean perturbations $\Delta g, \Delta f$ with
    covariances $\sfSigma{gg}$, $\sfSigma{ff}$ and
    cross-covariance $\sfSigma{gf}$, the second-order
    correction to the mean of the composed element is:
    \beq
    \delta\mu^{(2)} = \frac{1}{2}
    \sum_{\alpha,\beta} \langle \delta g_\alpha \, \delta f_\beta \rangle\,
    C_{\alpha\beta},
    \label{eq:meancorrection}
    \eeq
    where $C_{\alpha\beta}$ are the structure constants
    of $\mathfrak{se}(3)$ evaluated at the base point,
    and the expectation is taken over the joint distribution.
\end{theorem}

For the covariance, the second-order correction involves
the Hessian of the BCH map. We derive explicit expressions
using the finite composition formula~\eqref{eq:addSO3finite}
rather than the infinite BCH series, yielding:

\begin{remark}[Truncation order]
    The use of second-order versus third-order BCH truncations
    is justified quantitatively in Appendix~\ref{app:BCH},
    Table~\ref{tab:bchorder}. The second-order truncation
    error scales as $\mathcal{O}(\sigma^3)$ and the third-order
    as $\mathcal{O}(\sigma^4)$, with a measured ratio of
    $\sim 1/\sigma$. For typical rotational uncertainties
    $\sigma \leq 0.1$~rad, the second-order truncation
    already achieves $10^{-4}$~rad accuracy.
\end{remark}

\begin{theorem}[Second-order covariance correction]
    \label{thm:covcorrection}
    The covariance of $\ln h$ receives a correction of order
    $\mathcal{O}(\sigma^4)$ from the quadratic nonlinearity
    of the BCH map on $SE(3)$:
    \beq
    \sfSigma{hh}^{(2)} = \sfSigma{hh}^{(1)} +
    \Delta\Sigma^{(2)},
    \eeq
    where $\sfSigma{hh}^{(1)}$ is the first-order
    result~\eqref{eq:firstordercomp}. The correction
    $\Delta\Sigma^{(2)}$ is expressed entirely in terms of
    the covariance blocks $\sfSigma{aa}$, $\sfSigma{bb}$,
    $\sfSigma{ab}$ of the two perturbations (rotational and
    translational), weighted by the structure constants
    $\varepsilon_{ijk}$ and the coupling Jacobian
    $\mathsf{J}_t$. For Gaussian perturbations, no
    higher-order moments are needed.
    The explicit form is given in Appendix~\ref{app:secondorder}.
\end{theorem}

\subsection{Propagation of Points Through Uncertain Poses}

Given a 3D landmark distribution $p(\mathbf{x}) \sim
    \mathcal{N}(\bar{\mathbf{x}}, \sfSigma{xx})$ transformed
by an uncertain pose $\hat{f} = (f_0, \sfSigma{ff_r})$,
the marginalized distribution is:
\beq
p(\mathbf{x}_{\hat{f}}) = \int_{\real^6}
p(\mathbf{x}_{f_0 \cdot \Delta f})\, p(\ln \Delta f)\,
d\mathbf{t}\, d\bfOmega.
\label{eq:pointmargin}
\eeq

At first order:
\beq
p(\mathbf{x}_{\hat{f}}) \approx
\mathcal{N}\!\left(\bar{\mathbf{x}}_{f_0},\;
\sfSigma{xx_f} + \mathsf{J}_\star \sfSigma{ff_r}
\mathsf{J}_\star^T\right),\quad
\mathsf{J}_\star = \mathsf{R}_{f_0}\mathsf{J}_\times(\mathbf{x}),
\label{eq:pointfirstorder}
\eeq
with $\mathsf{J}_\times(\mathbf{x}) = [-[\mathbf{x}]_\times,\,
    \mathsf{I}]_{3 \times 6}$.

The second-order correction captures the non-Gaussian
character of~\eqref{eq:pointmargin}:

\begin{proposition}[Second-order point propagation]
    \label{prop:pointprop}
    The mean of $\mathbf{x}_{\hat{f}}$ receives a
    second-order correction
    \beq
    \langle \mathbf{x}_{\hat{f}} \rangle =
    \bar{\mathbf{x}}_{f_0} +
    \frac{1}{2}\mathrm{tr}\!\left(
    \frac{\partial^2 (f \star \mathbf{x})}{\partial (\ln \Delta f)^2}
    \bigg|_{f_0} \sfSigma{ff_r}\right)
    + \mathcal{O}(\sigma^3),
    \eeq
    where the Hessian $\partial^2(f\star\mathbf{x})/\partial(\ln\Delta f)^2$
    is computed from the $SE(3)$ action and involves the
    rotation-translation coupling.
\end{proposition}

%% ====================================================================
\section{Saddlepoint Marginalization}
\label{sec:saddlepoint}
%% ====================================================================

\subsection{The Marginalization Problem}

Consider the joint posterior for pose $f$ and 3D
landmarks $\{\mathbf{x}^j\}$ given projective
observations $\mathbf{z}$:
\beq
p(f, \{\mathbf{x}\} | \mathbf{z}) \propto
\prod_j p(\mathbf{z}^j | f, \mathbf{x}^j)\,
p(\mathbf{x}^j)\, p(f),
\label{eq:jointposterior}
\eeq
with the projective likelihood
\beq
p(\mathbf{z} | f, \mathbf{x}) \propto
\exp\!\left[-\frac{1}{2}
(\mathbf{z} - \pi(f \star \mathbf{x}))^T
\sfSigma{zz}^{-1}
(\mathbf{z} - \pi(f \star \mathbf{x}))\right],
\label{eq:projlikelihood}
\eeq
where $\pi(\mathbf{x}') = [x_1'/x_3', x_2'/x_3', 1]$
is the projective function.

The product rule gives
\beq
p(f, \{\mathbf{x}\} | \mathbf{z}) =
p(f | \mathbf{z}) \prod_j p(\mathbf{x}^j | f, \mathbf{z}).
\label{eq:productrule}
\eeq

\subsection{Standard Approach: Laplace Approximation (Schur Complement)}

The standard approach in bundle adjustment computes
the MAP estimate jointly, or uses the Schur complement
to eliminate landmarks---equivalent to a Laplace
(Gaussian) approximation of each $p(\mathbf{x}^j | f, \mathbf{z})$.

However, because $\pi$ is a rational function of $f \star \mathbf{x}$,
the conditional $p(\mathbf{x}^j | f, \mathbf{z})$ can exhibit
significant skewness and kurtosis, particularly when:
(a)~the landmark is near the optical axis ($x_3' \to 0$),
(b)~few observations constrain the landmark, or
(c)~the measurement noise is not small relative to the
depth $x_3'$.

\subsection{The Saddlepoint Approximation}

We seek the marginal
\beq
p(f | \mathbf{z}) = \prod_j \int
p(f, \mathbf{x}^j_{\rm opt} | \mathbf{z})\,
d\mathbf{x}^j_{\rm opt},
\label{eq:marginalize}
\eeq
via a two-step scheme:

\textbf{Step 1} (Optimize landmarks): For fixed $f$,
find $\mathbf{x}^j_{\rm opt}(f)$ by solving the
(linear) landmark optimization.

\textbf{Step 2} (Saddlepoint marginalization):
Approximate each integral in~\eqref{eq:marginalize}
using the saddlepoint method.

The saddlepoint approximation to the marginal
contribution of landmark $j$ is:
\beq
\int p(f, \mathbf{x}^j | \mathbf{z})\, d\mathbf{x}^j
\approx
p(f, \mathbf{x}^j_{\rm opt} | \mathbf{z})\,
(2\pi)^{3/2}\, |\mathsf{H}_j(f)|^{-1/2}\,
\exp\!\left(\frac{1}{6}\sum_{abc}
\kappa_{abc}^{(j)} s_a s_b s_c\right),
\label{eq:saddlepoint}
\eeq
where $\mathsf{H}_j(f)$ is the Hessian of the
negative log-posterior with respect to $\mathbf{x}^j$
at $\mathbf{x}^j_{\rm opt}$, $\kappa_{abc}^{(j)}$
are the third cumulants, and $s_a$ is the saddlepoint
solution.

\begin{remark}
    The Laplace approximation is the special case where
    the third cumulant correction is dropped. The saddlepoint
    correction captures the asymmetry introduced by the
    projective function $\pi$ and is $\mathcal{O}(1/n)$
    more accurate than Laplace for $n$ observations.
\end{remark}

\subsection{Derivatives for the Saddlepoint Correction}

The projective observation model~\eqref{eq:projlikelihood}
linearized at $\mathbf{x}'_t = f_t \star \mathbf{x}_t$ gives:
\beq
\mathsf{P} = \frac{1}{x_3'}\begin{pmatrix}
    1 & 0 & -u \\ 0 & 1 & -v
\end{pmatrix},\quad
\pi(\mathbf{x}') = [u, v, 1].
\label{eq:projjacobian}
\eeq

The measurement information matrix
for landmark $j$ observed from camera $i$ is:
\beq
\mathsf{M}(f_i, \mathbf{x}^j) =
\mathsf{R}_{f_i}^T \mathsf{P}_t^T\,
\sfSigma{zz}^{-1}\, \mathsf{P}_t\, \mathsf{R}_{f_i}.
\label{eq:Mmatrix}
\eeq
The full measurement information matrix is block-sparse, with nonzero off-diagonal blocks only where landmark $j$ is visible from camera $i$; the saddlepoint marginalization exploits this structure by integrating each landmark's $3\times 3$ block independently, reducing the system from $(6n+3m)$ to $6n$ dimensions.

The third-order derivatives required for~\eqref{eq:saddlepoint}
involve $\partial^3 \pi / \partial \mathbf{x}'^3$, which
are rational functions of $x_3'$ and can be computed in
closed form from the projective model.

\textbf{Step 3} (Pose optimization on $SE(3)$):
With the saddlepoint-corrected marginal
$p(f | \mathbf{z})$ in hand, we apply Gauss-Newton
iteration on $SE(3)$ using the compositional update
\beq
f_{t+1} = f_t \cdot \exp(\Delta \mathbf{f}_t),
\label{eq:poseupdate}
\eeq
reducing the problem to a $6n$-dimensional system
for $n$ camera poses. The required Jacobian of the
observation with respect to a right deviation is:
\beq
\frac{\partial(f \star \mathbf{x})}{\partial \delta_r f} =
\mathsf{R}_f\, \mathsf{J}_\times(\mathbf{x}),\quad
\mathsf{J}_\times(\mathbf{x}) =
[-[\mathbf{x}]_\times,\, \mathsf{I}]_{3 \times 6}.
\label{eq:posejacobian}
\eeq

\subsection{Comparison: Saddlepoint vs.\ Laplace vs.\ Joint MAP}

\begin{itemize}
    \item \textbf{Joint MAP} (bundle adjustment):
          Optimizes over $(6n + 3m)$ variables simultaneously.
          Accurate but computationally expensive ($\mathcal{O}((6n+3m)^2)$
          for dense problems) and sensitive to initialization
          due to non-convexity.

    \item \textbf{Laplace marginalization} (Schur complement):
          Marginalizes landmarks assuming Gaussian conditional.
          Reduces to $6n$ dimensions. Equivalent to standard
          BA with Schur complement.

    \item \textbf{Saddlepoint marginalization} (this work):
          Same $6n$-dimensional reduction as Laplace, but with
          a correction term capturing non-Gaussian effects.
          Larger convergence basin due to better handling of
          the projective non-convexity.
\end{itemize}

%% ====================================================================
\section{Numerical Experiments}
\label{sec:experiments}
%% ====================================================================

\subsection{Experiment 1: Bias in Parametrization}

We verify Proposition~\ref{prop:bias} using a
3D point cloud localization scenario. $N = 3$ landmarks are
randomly distributed in the cube $[-1,1]^3$ around the origin.
The camera is located at $\mathbf{T}_{f_0} = [0,0,-10]$ with
its view directed towards the center. We corrupt the original
landmarks with additive Gaussian noise of zero mean and
variance $\sigma = 1.0$, then solve the localization
problem~\eqref{eq:localization} using Horn's quaternion
method~\cite{Horn87} for each noise realization.
The resulting sample distribution of estimated camera poses
is plotted in two coordinate systems.

\begin{figure}[t]
    \centering
    \includegraphics[width=4.2cm]{\getname{SE3addomega}}
    \includegraphics[width=4.2cm]{\getname{SE3addT}}
    \caption{Projected sample distribution in second-kind coordinates:
        rotational components $\Delta\bfOmega = \bfOmega - \bfOmega_{f_0}$ (left)
        and translational coordinates $\Delta\mathbf{T} = \mathbf{T} - \mathbf{T}_{f_0}$ (right).
        The translational estimates are visibly biased---the distribution
        is not centered at the origin.}
    \label{fig:addanalysis}
\end{figure}

\begin{figure}[t]
    \centering
    \includegraphics[width=4.2cm]{\getname{SE3compomega}}
    \includegraphics[width=4.2cm]{\getname{SE3compt}}
    \caption{The same sample distribution as in Fig.~\ref{fig:addanalysis},
    now plotted in $SE(3)$ exponential coordinates (Lie-Cartan first kind),
    given by $\ln(f_0^{-1} \cdot f) = [\Delta\tilde{\bfOmega}, \mathbf{t}]$.
    The distribution is centered around the origin and exhibits
    approximately Gaussian behavior, confirming Proposition~\ref{prop:bias}.}
    \label{fig:companalysis}
\end{figure}

Fig.~\ref{fig:addanalysis} shows the sample distribution in
second-kind coordinates $[\bfOmega, \mathbf{T}]_{L_2}$. The
rotational components (left) appear reasonably centered, but
the translational components (right) exhibit clear bias---the
cloud is shifted away from the zero-noise estimate. This is
the bias mechanism of Proposition~\ref{prop:bias}: the nonlinear
coupling $\mathbf{T} = \mathsf{S}(\bfOmega)\mathbf{t}$
together with nonvanishing cross-correlations
$\langle \Omega_\alpha t_\beta \rangle \neq 0$ produces a
nonzero mean for $\mathbf{T}$ even when $[\bfOmega, \mathbf{t}]$
is unbiased.

Fig.~\ref{fig:companalysis} shows the identical sample set
plotted in exponential coordinates $[\bfOmega, \mathbf{t}]_{L_1}$.
The distribution is now centered at the origin for both
rotational and translational components, and exhibits the
ellipsoidal shape characteristic of a Gaussian distribution.
This empirically validates the choice of Lie-Cartan first-kind
coordinates for unbiased parametric uncertainty
representation~\eqref{eq:gaussianSE3}.

\subsection{Experiment 2: Second-Order Propagation Accuracy}

We validate the second-order mean and covariance corrections
(Theorems~\ref{thm:meanshift} and~\ref{thm:covcorrection})
against Monte Carlo ground truth.
Two zero-mean Gaussian perturbations
$\delta g \sim \mathcal{N}(0, \sfSigma{aa})$ and
$\delta f \sim \mathcal{N}(0, \sfSigma{bb})$
with cross-covariance $\sfSigma{ab} \neq 0$ are drawn
and composed via the second-order BCH map
$\bfOmega_c = \bfOmega_a + \bfOmega_b
    + \tfrac{1}{2}\bfOmega_a \times \bfOmega_b$.
We use $10^5$ samples and covariances with
$\sigma_\omega \approx 0.2$~rad, representative of
moderate rotational uncertainty.

\paragraph{Mean correction.}
The analytic second-order mean shift (Theorem~\ref{thm:meanshift})
predicts a correction
$\delta\mu^{(2)} = [-0.50, -1.00, -2.50] \times 10^{-3}$~rad,
arising from the antisymmetric part of $\sfSigma{ab}$
contracted with the structure constants~$\varepsilon_{i\alpha\beta}$.
The Monte Carlo mean is
$\langle \bfOmega_c \rangle_{\rm MC}
    = [-0.44, -1.66, -2.82] \times 10^{-3}$~rad.
The agreement in direction and magnitude confirms
the formula; the residual ($7.4\times 10^{-4}$~rad)
is consistent with the $\mathcal{O}(\sigma^3)$ truncation
of the BCH series (cf.\ Table~\ref{tab:bchorder}
at $s \approx 0.2$).

\paragraph{Covariance correction.}
The Isserlis-based formula~\eqref{eq:covcorrection_explicit}
yields a symmetric correction matrix
$\Delta\Sigma^{(2)}$ with dominant diagonal entries
of order $10^{-3}$.
Table~\ref{tab:covmc} compares the first-order and
corrected covariances against the Monte Carlo reference
in terms of the Frobenius norm of the error.

\begin{table}[t]
    \centering
    \caption{Covariance propagation accuracy for the rotational
        block ($3\times 3$) at $\sigma_\omega \approx 0.2$\,rad,
        measured in Frobenius norm of the error relative to
        $10^5$-sample Monte Carlo reference.}
    \label{tab:covmc}
    \begin{tabular}{l|c}
        Method                                         & $\|\Sigma - \Sigma_{\rm MC}\|_F$ \\
        \hline
        First-order~\eqref{eq:firstordercomp}          & $2.19 \times 10^{-3}$            \\
        Second-order~\eqref{eq:covcorrection_explicit} & $1.20 \times 10^{-3}$            \\
        \hline
        Improvement factor                             & $1.83\times$
    \end{tabular}
\end{table}

The second-order correction reduces the covariance error
by a factor of~$1.83$.
While the absolute improvement is modest at this noise level---the
correction $\Delta\Sigma^{(2)} \sim 10^{-3}$ is a perturbation
on the first-order covariance $\Sigma^{(1)} \sim 10^{-1}$---three
observations argue for its practical relevance:
\begin{enumerate}
    \item The correction scales as $\sigma_\omega^4$
          (Theorem~\ref{thm:covcorrection}), so the improvement
          factor grows quadratically with uncertainty.
          At $\sigma_\omega = 0.5$~rad the correction
          exceeds $10\%$ of the first-order covariance.
    \item In sequential estimation (EKF, preintegration),
          composition is applied at every time step.
          A systematic $\sim\!2\times$ per-step improvement
          compounds over long trajectories.
    \item The correction has \emph{no additional computational cost}
          beyond evaluating the structure constants and
          covariance blocks that are already available
          in the filter state---it is a closed-form formula,
          not a sampling procedure.
\end{enumerate}

\subsection{Experiment 3: Marginalized MAP vs.\ Bundle Adjustment}

We test the marginalized MAP algorithm on natural terrain
stereo image pairs. Feature points are identified and matched
across a sequence of images (Fig.~\ref{fig:matchedpoints}).
Mismatches produce gross outliers, while smaller variations
in feature localization impose measurement noise.

\begin{figure}[t]
    \centering
    \includegraphics[width=7cm]{\getname{realimage}}
    \caption{Natural terrain scenario with matched feature points.
        Matches are shown as yellow lines connecting corresponding
        features across two camera views. Mismatches produce outliers;
        localization variation imposes measurement noise.}
    \label{fig:matchedpoints}
\end{figure}

\begin{figure}[t]
    \centering
    \includegraphics[width=4.2cm]{\getname{costplot1}}
    \includegraphics[width=4.2cm]{\getname{costplot2}}
    \caption{Convergence comparison: log-likelihood (log scale) vs.\
        iteration steps. Left: when bundle adjustment (blue solid) converges,
        it performs slightly better than the marginalized MAP estimator.
        Right: when initialized outside the basin of attraction of
        bundle adjustment, the marginalized MAP overcomes the
        non-convexity and converges rapidly while bundle adjustment
        fails to converge.}
    \label{fig:performance}
\end{figure}

Fig.~\ref{fig:performance} compares the convergence behavior.
Both methods converge within approximately 10 iterations when
the initialization lies within the basin of attraction (left panel).
However, the convergence region of the marginalized MAP algorithm
is considerably larger (right panel), since the non-convexity
of the projective cost function is handled by the
marginalization step~\eqref{eq:marginalize}. The joint estimate
requires operations on a $210 \times 210$ matrix (68 landmarks,
6 DOF for camera), whereas the marginalized MAP operates on
$6 \times 6$ matrices.

\subsection{Experiment 4: Pose Estimate Distribution Quality}

To assess the statistical quality of pose estimates, we impose
Gaussian noise ($\sigma = 0.8\%$ of image dimensions,
corresponding to $\pm 2$ pixel error for a $512$-pixel image)
on the feature point observations and compute the pose estimate
distribution. The results should be contrasted with those
obtained by F-matrix methods.

\begin{figure}[t]
    \centering
    \includegraphics[width=4.2cm]{\getname{LMedSqomega}}
    \includegraphics[width=4.2cm]{\getname{LMedSqtexp}}
    \caption{Sample distribution (2000 samples) of Least Median
        Squares estimates in exponential coordinates: rotation (left)
        and translation (right), from 68 data points corrupted with
        Gaussian noise $\sigma = 0.8\%$ of image dimensions.
        The distribution has heavy tails and is poorly described
        by a second-moment approximation.}
    \label{fig:LMedSq}
\end{figure}

\begin{figure}[t]
    \centering
    \includegraphics[width=4.2cm]{\getname{Bayesomega05}}
    \includegraphics[width=4.2cm]{\getname{Bayestexp05}}
    \caption{Sample distribution of camera pose estimates from
        the marginalized MAP algorithm in exponential coordinates:
        rotation (left) and translation (right), using the same
        matched data and noise level as Fig.~\ref{fig:LMedSq}.
        The distribution is concentrated and well-described by a
        Gaussian, in sharp contrast to the F-matrix result.}
    \label{fig:bayesnoise}
\end{figure}

Fig.~\ref{fig:LMedSq} shows the estimator distribution obtained
by the LMedS F-matrix method with reweighted least squares
refinement. Despite achieving some robustness, significant
samples remain far from the zero-noise estimate. The
distribution has heavy tails and is not adequately described
by a second-moment (Gaussian) approximation---precisely the
regime where the saddlepoint correction
of Section~\ref{sec:saddlepoint} becomes relevant.

Fig.~\ref{fig:bayesnoise} shows the corresponding distribution
from the marginalized MAP algorithm. The pose estimates are
tightly concentrated and exhibit the ellipsoidal shape of a
Gaussian distribution in exponential coordinates, confirming
that the Bayesian framework with proper $SE(3)$
parametrization yields well-behaved estimator statistics.

The contrast between Figs.~\ref{fig:LMedSq}
and~\ref{fig:bayesnoise} demonstrates two points simultaneously:
(i)~the marginalized MAP estimator produces dramatically better
pose statistics than algebraic (F-matrix) methods under the
same noise conditions, and (ii)~the exponential coordinate
parametrization yields distributions that are well-approximated
by a Gaussian, validating the uncertainty
model~\eqref{eq:gaussianSE3}.

%% ====================================================================
\section{Conclusion}
\label{sec:conclusion}
%% ====================================================================

We have presented three contributions to Bayesian
inference on $SE(3)$:

\begin{enumerate}
    \item A complete, closed-form algebra for $SE(3)$
          uncertainty propagation, including the coupling Jacobian
          $\mathsf{J}_t(\bfOmega, \mathbf{t})$, with explicit
          second-order corrections to both mean and covariance.

    \item A saddlepoint marginalization scheme that
          eliminates latent Euclidean variables from the joint
          posterior more accurately than the standard Laplace
          (Gauss-Newton) approximation.

    \item A closed-form BCH formula for finite Rodrigues
          vector composition via $SU(2)$, with careful treatment
          of phase reflection at $\Theta = \pi$.
\end{enumerate}

The framework is applicable beyond pose estimation
to any domain requiring probabilistic inference on
$SE(3)$: IMU preintegration, spacecraft attitude
determination, surgical robot registration, and
multi-body dynamics.

Future work will address: (i)~recursive (filtering)
implementations with the second-order corrections;
(ii)~extension to outlier-robust estimation via
MAPSAC-type methods; (iii)~application to
visual-inertial odometry with preintegration on
manifolds.

%% ====================================================================
\appendices
%% ====================================================================

\section{$SE(3)$ Matrix Group}
\label{app:SE3group}
A group element $g\in SE(3)$ can be represented by a real valued $4\times 4$ matrix
\beq
g \cong (\mathsf{R}(\bfOmega_g),{\bf T}_g) \equiv \left(
\begin{array}{cc}
        \mathsf{R}(\bfOmega_g) & {\bf T}_g \\
        0\; 0\; 0              & 1
    \end{array} \right) \label{eq:SE3matrixrepresentation}
\eeq
and group operation is endowed as usual by matrix multiplication.
Throughout this appendix, $[{\bf v}]_\times$ denotes the skew-symmetric matrix associated
with a vector ${\bf v}\in\real^3$, defined by $[{\bf v}]_\times\,{\bf w} = {\bf v}\times{\bf w}$
for all ${\bf w}\in\real^3$.

\subsection{Exponential and logarithmic map}
\label{app:expmap}
Any matrix of the structure (\ref{eq:SE3matrixrepresentation}) can be obtained
by the exponential map
\beq
\exp: \real^6 \rightarrow SE(3),\;
[\bfOmega_g, {\bf t}_g]_{L_1} \rightarrow
\exp( \hat{\bf H}\cdot\bfOmega_g + \hat{\bf T}\cdot{\bf t}_g ) = g,
\label{eq:exponentialmap}
\eeq
where the coordinates $[\bfOmega_g,{\bf t}_g]_{L_1} \in \real^6$ are called the
    {\em Lie-Cartan coordinates of the first kind}. We usually will omit the notation subscript
$L_1$ when dealing with those {\em exponential} coordinates. The exponential map is globally well defined,
and can be depicted geometrically as the path connecting the identity element $1\in SE(3)$ with
$g$. In (\ref{eq:exponentialmap}), the basis elements
of the 6-dimensional vector space $\real^6\cong T_1SE(3)$ are compactly denoted as,
\beq
\hat{\bf H}\cdot\bfomega=\mathsf{H}_x\Omega_x+\mathsf{H}_y\Omega_y+\mathsf{H}_z\Omega_z,\quad
\hat{\bf T}\cdot{\bf t} =\mathsf{T}_x t_x + \mathsf{T}_y t_y + \mathsf{T}_z t_z,
\label{eq:generatingmatrices}
\eeq
where each element of $\hat{\bf H}$ and $\hat{\bf T}$ are real $4\times 4$ matrices,
also called the generators of the $SE(3)$ group. In conjunction with the Lie-product
({\em commutator}), $[\mathsf{A},\mathsf{B}]_c = \mathsf{AB}-\mathsf{BA}$\footnote{The square bracket
    notation is frequently used without the subscript $c$; we use it here because we already
    defined $[{\bfomega},{\bf v}]$ as the concatenation of 2 vectors.},
the vector space becomes the $\mathfrak{se}(3)$ Lie algebra,
\beq
[\mathsf{H}_i,\mathsf{H}_j]_c = \varepsilon_{ijk}\mathsf{H}_k,\;
[\mathsf{H}_i,\mathsf{T}_j]_c = \varepsilon_{ijk}\mathsf{T}_k,\;
[\mathsf{T}_i, \mathsf{T}_j]_c = 0. \label{eq:commutatortable}
\eeq
The commutator can be seen as the adjoint operator,
\beq
\adj{\mathsf{A}} \mathsf{B} = [\mathsf{A},\mathsf{B}]_c, \;\;
\adj{\mathsf{A}}^k(\mathsf{B}) = [\mathsf{A},\ldots,[\mathsf{A},[\mathsf{A},\mathsf{B}]_c]_c,
\label{eq:defadjoint}
\eeq
which gives a powerful notation to represent $k$ interlocked commutators.

Frequently, we will have to deal with a factorized representation of (\ref{eq:exponentialmap})
\begin{multline}
    ( \mathsf{R}_z(\phi)\mathsf{R}_y(\vartheta)\mathsf{R}_z(\psi),{\bf T}) = \exp( \hat{\bf T}\cdot {\bf T} ) \\
    \exp( \mathsf{H}_z\phi)\exp( \mathsf{H}_y\vartheta)\exp( \mathsf{H}_z\psi),
    \label{eq:Lie2ndkind}
\end{multline}
where $({\bf T},\phi,\vartheta,\psi)_{L_2}$ are called {\em canonical coordinates of second kind}\footnote{The ordering of the exponential factors is important, hence there are multiple sets
    of coordinates, i.e.\ Euler ZYZ, XYZ$\ldots$}.
Due to the non-trivial algebra structure (\ref{eq:commutatortable}) the connection between first
and second kind coordinates is not a simple one.
However, we can give a simple commutation rule for the following case
\beq
\exp(\hat{\bf T}\cdot{\bf T})\exp(\hat{\bf H}\cdot\bfOmega) =
\exp(\hat{\bf H}\cdot\bfOmega)\exp(\hat{\bf T}\cdot\mathsf{R}^{-1}(\bfOmega){\bf T}),
\label{eq:TOmegacommute}
\eeq
which can be derived from the general conjugation rule
\beq
g^{-1}\,\exp(\hat{\bf H}\cdot\bfOmega+\hat{\bf T}\cdot{\bf T})\, g =
\exp(g^{-1}\; (\hat{\bf H}\cdot\bfOmega+\hat{\bf T}\cdot{\bf T})\; g) \label{eq:conjugation}
\eeq
where $g$ in the exponent of (\ref{eq:conjugation}) is interpreted as a $4\times 4$ matrix, according to
(\ref{eq:SE3matrixrepresentation}).

In particular, for the exponential map (\ref{eq:exponentialmap}) we can give the following Eqns.
\beqar
\mathsf{R}(\bfOmega) &=& \mathsf{I}+\sin\Theta\, \mathsf{H} + (1-\cos\Theta)\mathsf{H}^2,
\quad \label{eq:rodriguesapp} \\
{\bf T} &=& \mathsf{S}(\bfOmega)\; {\bf t},\quad
\mathsf{S}(\bfOmega) = \mathsf{I} + \mathsf{H}^2 +  {1\over\Theta}(\mathsf{I}-\mathsf{R})\mathsf{H}, \nonumber
\eeqar
where we use the notation
\beq
\Theta=|\bfOmega|,\; {\bf r} = \bfOmega/\Theta,\quad
\mathsf{H}=\hat{\bf H}\cdot{\bf r}. \label{eq:expnotation}
\eeq
As seen from the Rodrigues formula (\ref{eq:rodriguesapp}), $\mathsf{R}(\bfOmega)$ is $2\pi$-periodic in $\Theta$.
This implies that a change in the normal vector direction can be compensated by a $2\pi$ offset,
$({\bf r},\Theta) \cong (-{\bf r},2\pi-\Theta)$. Hence, $\bfOmega$ uniquely specifies a rotation
only when it is restricted within a sphere of radius $\pi$. Penetrating the skin of this sphere
causes {\em phase reflection}, where opposite points on the boundary are identified,
$({\bf r},\pi) \equiv (-{\bf r},\pi)$.
This phenomenon has a precise topological origin in the double-cover structure
$SU(2)/\mathbb{Z}_2 \cong SO(3)$, which we develop in the next subsection.

The inverse function to the exponential map is the logarithmic map
\beq
\ln: SE(3)\rightarrow\real^6,\; (\mathsf{R}(\bfOmega),{\bf T})\rightarrow [\bfOmega,{\bf t}]_{L_1},
\label{eq:logarithmicmap}
\eeq
which is well defined for $\Theta<\pi$. As discussed above, the logarithmic map will not be
unique at the phase reflection boundary, $\Theta = \pi$. The existence of such boundary will have
implications for the topics considered in the following sections.

%% -----------------------------------------------------------------------
\subsection{The $SU(2)/\mathbb{Z}_2 \cong SO(3)$ connection}
\label{app:su2so3}
%% -----------------------------------------------------------------------
The $SO(3)$ generating algebra $\mathsf{H}_i$ and the $SU(2)$ one with generators $\mathsf{u}_i$
have very similar structures,
meaning there is an isomorphism between the algebras, $\mathfrak{so}(3) \cong \mathfrak{su}(2)$,
\beq
[\mathsf{H}_i,\mathsf{H}_j]_c = \varepsilon_{ijk}\mathsf{H}_k,\;
[\mathsf{u}_i,\mathsf{u}_j]_c = 2\varepsilon_{ijk}\mathsf{u}_k\;. \label{eq:su2-algebra-structure}
\eeq
Hence, after the rescaling $\tau_i = \tfrac{1}{2}\mathsf{u}_i$, we obtain
$[\tau_i,\tau_j]_c = \varepsilon_{ijk}\tau_k$, so that the correspondence
$\mathsf{H}_i \leftrightarrow \tau_i$ defines the Lie-algebra isomorphism
$\mathfrak{so}(3) \cong \mathfrak{su}(2)$.

Both $\mathsf{H}_0$ and $\mathsf{u}_0$ are the identity matrices in each of their cases;
they are the identity elements of the group representations, not generators
of the Lie algebra.
The standard matrix representations are:
\beq
\mathsf{u}_1 = \begin{pmatrix} 0 & i \\ i & 0 \end{pmatrix}, \;
\mathsf{u}_2 = \begin{pmatrix} 0 & -1 \\ 1 & 0 \end{pmatrix}, \;
\mathsf{u}_3 = \begin{pmatrix} i & 0 \\ 0 & -i \end{pmatrix}
\eeq
and for the Lie algebra $\mathfrak{so}(3)$ generators of the $SO(3)$ group
\begin{IEEEeqnarray*}{c}
    \mathsf{H}_1 = \begin{pmatrix} 0 & 0 & 0 \\ 0 & 0 & -1 \\ 0 & 1 & 0 \end{pmatrix},\quad
    \mathsf{H}_2 = \begin{pmatrix} 0 & 0 & 1 \\ 0 & 0 & 0 \\ -1 & 0 & 0 \end{pmatrix},\\
    \mathsf{H}_3 = \begin{pmatrix} 0 & -1 & 0 \\ 1 & 0 & 0 \\ 0 & 0 & 0 \end{pmatrix}\, .
\end{IEEEeqnarray*}
Other representations of $\mathfrak{su}(2)$ are possible and standard in Physics, with the
algebra structure (\ref{eq:su2-algebra-structure}) changing slightly to
$[{\sigma}_i,{\sigma}_j]_c = 2i\varepsilon_{ijk}{\sigma}_k$,
which introduces the Pauli matrices:
\beq
\sigma_1 = -i \mathsf{u}_1\,,\quad \sigma_2 = i \mathsf{u}_2\,,\quad \sigma_3 = -i \mathsf{u}_3\, .
\eeq

Regardless of the isomorphic $\mathfrak{su}(2)$ representations, the generated group
$SU(2)$ is not isomorphic to $SO(3)$: $SU(2)$ is the simply connected double cover
of $SO(3)$. To see this explicitly through the Rodrigues formulas, let
\beq
\mathsf{u} = \hat{\mathsf{u}}\cdot{\bf r} = r_i\,\mathsf{u}_i\,,\qquad
\Theta = |\bfOmega|\,,\qquad
{\bf r} = \bfOmega/\Theta\,.
\eeq
Since $\mathsf{u}^2 = -\mathsf{u}_0$ for any unit vector ${\bf r}$, the matrix exponential
in $SU(2)$ truncates to
\beq
U(\bfOmega) = \exp\!\Bigl(\tfrac{1}{2}\,\hat{\mathsf{u}}\cdot\bfOmega\Bigr)
= \mathsf{u}_0\cos\frac{\Theta}{2}
+ \mathsf{u}\,\sin\frac{\Theta}{2}\,. \label{eq:rodrigues-su2}
\eeq
The adjoint action of $U$ on the pure-imaginary quaternion subspace induces
a rotation $\mathsf{R}(\bfOmega)\in SO(3)$:
\beq
U(\bfOmega)\;(\hat{\mathsf{u}}\cdot{\bf x})\;U(\bfOmega)^{-1}
= \hat{\mathsf{u}}\cdot\bigl(\mathsf{R}(\bfOmega)\,{\bf x}\bigr)\,,
\eeq
\begin{multline}
    \mathsf{R}(\bfOmega)\,{\bf x}
    = \cos\Theta\;{\bf x}
    + (1-\cos\Theta)({\bf r}\cdot{\bf x})\,{\bf r} \\
    + \sin\Theta\; {\bf r}\times{\bf x}\,. \label{eq:rodrigues-vector}
\end{multline}
In matrix form this is precisely the Rodrigues formula (\ref{eq:rodriguesapp}).

Equivalently, writing $U$ in quaternion parametrisation,
\begin{gather}
    U = q_0\,\mathsf{u}_0 + q_i\,\mathsf{u}_i\,,\qquad
    q_0^2 + {\bf q}^{\!\top}{\bf q} = 1\,, \nonumber\\
    q_0 = \cos\frac{\Theta}{2}\,,\qquad
    {\bf q} = \sin\frac{\Theta}{2}\;{\bf r}\,,
\end{gather}
the induced rotation matrix becomes
\beq
\mathsf{R}(U) = (q_0^2 - {\bf q}^{\!\top}{\bf q})\,\mathsf{H}_0
+ 2\,{\bf q}\,{\bf q}^{\!\top}
+ 2\,q_0\,[{\bf q}]_\times\,. \label{eq:quaternion-to-R}
\eeq
This exhibits $SU(2)\cong S^3$ directly.

\paragraph{The covering map.}
The surjective homomorphism $\pi\colon SU(2)\to SO(3)$ is therefore given by
\beq
\pi(U) = \mathsf{R}(U)\,,\qquad
\ker\pi = \{\pm \mathsf{u}_0\}\,.
\eeq
Hence
\beq
SO(3) \cong SU(2)/\{\pm \mathsf{u}_0\} \cong S^3/\mathbb{Z}_2\,.
\eeq

\paragraph{The two-to-one identification.}
In the exponential parametrisation the double cover appears as
\beq
U(\bfOmega + 2\pi{\bf r}) = -\,U(\bfOmega)\,,\qquad
\pi(U) = \pi(-U)\,.
\eeq
Antipodal points in $SU(2)\cong S^3$ therefore represent the same rotation
in $SO(3)$. This is the precise topological origin of the
phase-reflection phenomenon encountered in the previous subsection:
the identification $({\bf r},\pi) \equiv (-{\bf r},\pi)$ on the boundary
of the Rodrigues parameter sphere is the quotient $S^3/\mathbb{Z}_2$
restricted to the equator. Equivalently,
$\pi_1(SO(3)) = \mathbb{Z}_2$, whereas $SU(2)\cong S^3$ is simply connected.

\paragraph{Inverse map.}
Given $\mathsf{R}\in SO(3)$, the rotation angle follows from $\mathrm{Tr}\,\mathsf{R} = 1 + 2\cos\Theta$
and the axis from the antisymmetric part,
$n_k = -\varepsilon_{ijk}\,R_{ij}/(2\sin\Theta)$, after which (\ref{eq:rodrigues-su2})
returns the $SU(2)$ element up to an overall sign.
Alternatively, reading off from (\ref{eq:quaternion-to-R}) directly:
$q_0 = \pm\tfrac{1}{2}\sqrt{1+\mathrm{Tr}\,\mathsf{R}}$ and
$q_k = -\varepsilon_{ijk}\,R_{ij}/(4q_0)$ for $q_0\neq 0$.
The sign ambiguity is precisely the $\mathbb{Z}_2$ kernel: no continuous global
section $SO(3)\to SU(2)$ exists.

%% ====================================================================
\section{Baker-Campbell-Hausdorff Formula}
\label{app:BCH}
%% ====================================================================
Frequently, we obtain a factorized rotation representation,
\beq
\exp( \Theta_c \mathsf{H}_c ) = \exp(\Theta_a \mathsf{H}_a)\exp(\Theta_b\mathsf{H}_b),
\label{eq:expproduct}
\eeq
presented with notation (\ref{eq:expnotation}). An exponential series expansion of the right hand side of
(\ref{eq:expproduct}) is ordered with $\mathsf{H}_a$ to the left and $\mathsf{H}_b$ to the right. It is clear
that a simple sum $\mathsf{H}_c \neq \mathsf{H}_a + \mathsf{H}_b$ cannot hold because its exponential series expansion is unordered. The Baker-Campbell-Hausdorff (BCH) formula\footnote{The BCH formula can be derived using
    the derivative of exponentials, the resulting integral with function $h(x)=\ln(x)/(x-1)$ is understood as a power series expansion in $x$.} is useful to derive an {\em additive} relationship,
\beqar
\Theta_c \mathsf{H}_c &=& \Theta_b \mathsf{H}_b + \int_0^{\Theta_a} h\left(
e^{s \, \adj{\mathsf{H}_a}} e^{ \Theta_b \, \adj{\mathsf{H}_b}}
\right) \, \mathsf{H}_a \; ds, \nonumber \\
&=& \Theta_a \mathsf{H}_a + \Theta_b \mathsf{H}_b + {1\over 2}\Theta_a\Theta_b
[\mathsf{H}_a,\mathsf{H}_b]+
\ldots\; . \label{eq:BCH}
\eeqar
However, (\ref{eq:BCH}) only hands an infinite series expansion. It is quite useful to have a closed expression for the Rodrigues vector $\bfOmega_c$ of the resulting rotation $\mathsf{R}_c = \exp( \hat{\bf H}\cdot\bfOmega_c )$. We could compactly denote this operation as $\bfOmega_c = \ln \mathsf{R}_a\mathsf{R}_b\,$. Using the quaternion
representation of $SU(2)$ developed in \S\ref{app:su2so3}, one can show that
the following Eqn.\ holds,
\begin{multline}
    \bfOmega_c = {\Theta_c\over C}\left(
    a \, {\bf r}_a + b \, {\bf r}_b + ab \, {\bf r}_a\times{\bf r}_b \right), \\
    a =\tan{\Theta_a\over 2},\, b=\tan{\Theta_b\over 2}.
    \label{eq:addSO3finite}
\end{multline}
The normalizing factor, $C = \sqrt{A - \varepsilon^2}$, can be given as
\beq
A = 1+a^2+b^2+a^2b^2,\;\; \varepsilon = (1-ab\cos\vartheta),\quad
\cos\vartheta = {\bf r}_a\cdot{\bf r}_b\;, \label{eq:normfactor}
\eeq
and the norm $\Theta_c$
\beq
\tan{\Theta_c\over 2} = C / \, \varepsilon\;.
\eeq
For the trivial case, $\cos\vartheta=\pm1$ in (\ref{eq:normfactor}), we simply obtain
$\bfOmega_c=\bfOmega_a+\bfOmega_b$.
To evaluate the prefactor $\Theta_c/C$ in (\ref{eq:addSO3finite}) or the primed versions
below, we have to pay attention to the following two limit cases:
\begin{enumerate}
    \item limit $\varepsilon \rightarrow 0$, (this implies $A \geq 1$)
          \beq
          {\Theta_c\over C} = {\pi \over \sqrt{A}}\, \left[ 1 - {2\over\pi}\,{\varepsilon\over\sqrt{A}} + {1\over 2}\,
          \left({\varepsilon\over \sqrt{A}}\right)^2  +
          {\cal O}((\varepsilon/\sqrt{A})^3)  \right]
          \eeq
    \item limit $C\rightarrow 0$, (this implies $\varepsilon^2\geq 1$)
          \beq
          {\Theta_c\over C} = {2\over \varepsilon}\, \left[1 - {1\over 3}\,\left({C\over\varepsilon}\right)^2  +
          {\cal O}((C/\varepsilon)^4)  \right] .
          \eeq
\end{enumerate}

\subsection{Phase reflection at $\Theta=\pi$}
Since the logarithmic map loses its uniqueness at $\Theta=\pi$
(cf.\ the double-cover discussion in \S\ref{app:su2so3}),
we have to be careful to define the
vector addition (\ref{eq:addSO3finite}). Firstly, we consider
$\Theta_a \rightarrow \pi-$ which implies $a\rightarrow \infty$ and $b$ finite.
Secondly a similar statement accounts for the limit $a\rightarrow\infty$ and $b\rightarrow\infty$.
We notice that those two limits, (i) and (ii), can be treated consistently.
In both cases a cutline at $\cos\vartheta=0$ has to be made to account for the phase
reflection boundary. The limit process $\Theta_a \rightarrow \pi-$ can be depicted graphically,
see Fig.~\ref{fig:pwrap}.
\begin{figure}[t]
    \centering
    \includegraphics[width=6cm]{\getname{phasewrap}}
    \caption{Phase reflection occurs for $|\bfOmega_a| = \pi$. The result of {\em adding} another
        infinitesimal vector ${\bfomega}_b$ depends whether
        it has the same direction $\cos\vartheta>0$, which will give a resulting
        vector in the direction $-\bfOmega_a$. If $\cos\vartheta<0$, the resulting vector will be
        unchanged.}
    \label{fig:pwrap}
\end{figure}
Hence, all quantities are properly defined in the limits. Therefore, Eqn.~(\ref{eq:addSO3finite}) can
be used to {\em add} finite Rodrigues vectors without employing intermediate rotation matrices or quaternions.

\subsection{Adding two infinitesimal rotations: kinematics of rigid bodies}
The rate of change, {\em angular velocity}, of $\bfOmega$ for an infinitesimal small quantity $dt\rightarrow 0$
is $d\bfOmega = dt\, \bfomega$. We can apply the finite {\em addition} formula (\ref{eq:addSO3finite}) or
directly apply the BCH expansion (\ref{eq:BCH}) to obtain,
\beq
d\bfOmega_c = dt\, \bfomega_c = dt\, \bfomega_a + dt\, \bfomega_b + {1\over 2} dt^2\, \bfomega_a\times\bfomega_b\;.
\label{eq:addSO3inf}
\eeq
As such we recover in (\ref{eq:addSO3inf}) the well known fact that angular velocities $\bfomega$
constitute an algebraic vector space with the normal commutative addition. More exactly, the
vector space is known as $T_1SO(3)$.

\subsection{Truncation order of the BCH series}
\label{app:bchorder}

The BCH series~(\ref{eq:BCH}) provides successive approximations to
the exact composition $\bfOmega_c = \ln \mathsf{R}_a\mathsf{R}_b\,$
given by~(\ref{eq:addSO3finite}). For uncertainty propagation
(Appendix~\ref{app:secondorder}), the relevant truncations are:
\begin{enumerate}
    \item \emph{Second-order:} $\bfOmega_c \approx \bfOmega_a + \bfOmega_b
              + \tfrac{1}{2}\bfOmega_a\times\bfOmega_b$
    \item \emph{Third-order:} adds
          $\tfrac{1}{12}(\bfOmega_a\!\times\!(\bfOmega_a\!\times\!\bfOmega_b)
              + \bfOmega_b\!\times\!(\bfOmega_b\!\times\!\bfOmega_a))$
\end{enumerate}

To quantify the truncation error, we evaluate both approximations
against the exact closed-form result~(\ref{eq:addSO3finite})
for $\bfOmega_a = s\,\mathbf{r}_a$ and $\bfOmega_b = s\,\mathbf{r}_b$
with non-collinear unit vectors
$\mathbf{r}_a,\mathbf{r}_b$ and varying scale~$s$.
Table~\ref{tab:bchorder} confirms the expected asymptotic orders:
the second-order error scales as $\mathcal{O}(s^3)$ and the third-order
error as $\mathcal{O}(s^4)$.
Their ratio grows as~$1/s$, reflecting the one-order gain
from the additional commutator term.

\begin{table}[ht]
    \centering
    \caption{BCH truncation error versus scale $s = |\bfOmega_a| = |\bfOmega_b|$
    for non-collinear rotation axes.
    The ratio confirms $\mathrm{err}_{2\mathrm{nd}}/\mathrm{err}_{3\mathrm{rd}} \sim 1/s$.}
    \label{tab:bchorder}
    \begin{tabular}{c|cc|c}
        $s$ [rad] & $\mathrm{err}_{2\mathrm{nd}}$ & $\mathrm{err}_{3\mathrm{rd}}$ & ratio \\
        \hline
        0.001     & $1.28\times10^{-10}$          & $1.29\times10^{-14}$          & 9955  \\
        0.003     & $3.46\times10^{-9}$           & $1.04\times10^{-12}$          & 3318  \\
        0.01      & $1.28\times10^{-7}$           & $1.29\times10^{-10}$          & 996   \\
        0.03      & $3.46\times10^{-6}$           & $1.04\times10^{-8}$           & 332   \\
        0.1       & $1.28\times10^{-4}$           & $1.29\times10^{-6}$           & 99.5  \\
        0.3       & $3.47\times10^{-3}$           & $1.06\times10^{-4}$           & 32.9  \\
        1.0       & $1.31\times10^{-1}$           & $1.44\times10^{-2}$           & 9.1   \\
    \end{tabular}
\end{table}

For Gaussian rotational uncertainties with standard deviation
$\sigma_\omega$, the second-order BCH approximation introduces a
truncation bias of order $\sigma_\omega^3$ in the mean correction
and $\sigma_\omega^6$ in the covariance correction.
At $\sigma_\omega = 0.1$\,rad ($\approx 5.7^\circ$), the truncation
error is below $10^{-4}$\,rad, well within the linearization error
of typical EKF implementations.
For larger uncertainties ($\sigma_\omega \gtrsim 0.3$\,rad),
the third-order term should be retained or the exact
closed-form~(\ref{eq:addSO3finite}) used directly.

%% ====================================================================
\section{Wei-Norman Formula and Composition Jacobians}
\label{app:weinorman}
%% ====================================================================

\subsection{Right Rodrigues Jacobian}
We now derive the right Jacobian $\Jom$ that relates infinitesimal
perturbations in Lie-algebra coordinates to the corresponding group-level variations.
We present two routes: first via the Wei-Norman formula using the derivative of
exponentials, then a direct computation from the finite addition formula
(\ref{eq:addSO3finite}).

We consider the rotation product expression in (\ref{eq:expproduct}). We
describe the infinitesimal environment at $\bfOmega_c = \bfOmega_a + \Delta\bfomega$
for a finite left rotation $\bfOmega_a$ in dependence of an infinitesimal right rotation $\bfOmega_b \rightarrow \bf 0$.
Then we have to analyze the following directional derivative
\beq
{d\over ds}\;\exp(\hat{\bf H}\cdot\bfOmega_a + s\;\hat{\bf H}\cdot\Delta\bfomega)
\; = \; \exp(\hat{\bf H}\cdot\bfOmega_a) \hat{\bf H}\cdot\mathsf{J}_{\omega_r}^{-1}(\bfOmega_a)
\Delta\bfomega,
\label{eq:directionalderiv}
\eeq
where the right composition Jacobian $\mathsf{J}_{\omega_r}(\bfOmega_a) \in \real^{3\times3}:
    T_{\mathsf{R}(\bfOmega_a)}SO(3)\rightarrow T_1SO(3)$ describes the
mapping from the tangent space at $\mathsf{R}(\bfOmega_a)$ to the one at identity.
In order to evaluate the left hand side of (\ref{eq:directionalderiv}), it is advantageous to
consider a ``matrix'' equation for $g\in SE(3)$ and general $\mathsf{A},\mathsf{B}\in \mathfrak{se}(3)$
\beq
{\partial \over \partial t}g(t,s) = (\mathsf{A} + s\mathsf{B})g(s,t)\; \rightarrow \;
g(s,t) = e^{t (\mathsf{A}+s\mathsf{B})}. \label{eq:SE3diff}
\eeq
With the help of the Dyson ordering operator ${\cal T}$
and the adjoint notation (\ref{eq:defadjoint}), it can be shown that the following holds
\beq
g(s,t) = e^{t \mathsf{A}} \; {\cal T}\exp\left( s\,
\int_0^t e^{-u\, \adj{\mathsf{A}}} \, du \; \mathsf{B} \right ) \; .
\label{eq:timeordering}
\eeq
Differentiating (\ref{eq:timeordering}) with respect to $s$, setting $t=1$,
then applying it to (\ref{eq:directionalderiv}) yields
\beq
u(\Theta_a\, \adj{\mathsf{H}_a} )
\hat{\bf H}\cdot\Delta\bfomega = \hat{\bf H}\cdot \mathsf{J}_{\omega_r}^{-1}(\bfOmega_a)
\Delta\bfomega, \quad u(x)=\int_0^1 e^{-s x} \; ds
\label{eq:weinorman}
\eeq
a relationship to obtain $\mathsf{J}_{\omega_r}(\bfOmega_a)$. In fact (\ref{eq:weinorman}) constitutes a
Wei-Norman formula~\cite{Wei-Norman}, which relates the differential changes of a product of
exponentials (\ref{eq:expproduct}) to the changes of one resulting exponential. However, it is not easy
to obtain a closed form expression from (\ref{eq:weinorman}).

\paragraph{Direct route via finite addition.}
It shall be noticed that with Eqn.~(\ref{eq:addSO3finite})
we directly obtain the right Rodrigues Jacobian $\mathsf{J}_{\omega_r}(\bfOmega_a)$ by,
\beq
\mathsf{J}_{\omega_r}(\bfOmega_a) =
\left.{\partial \bfOmega_c \over \partial \bfOmega_b}\right|_{\Theta_b=0}=
\left.{\partial \bfOmega_c \over \partial \Theta_b}\right|_{\Theta_b=0} {\bf r}_b+
\left.{\partial \bfOmega_c \over \partial {\bf r}_b} {\partial {\bf r}_b \over \partial \bfOmega_b}
\right|_{\Theta_b=0}. \label{eq:jacobianomega1}
\eeq
Further, Eqn.~(\ref{eq:jacobianomega1}) can be evaluated as,
\beq
\mathsf{J}_{\omega_r}(\bfOmega_a) = (1-x)\, \mathsf{I}+ x\, {\bf r}_a{\bf r}_a^T + {1\over 2}[\bfOmega_a]_\times,
\; 1-x = {\Theta_a\over 2}\tan^{-1}{\Theta_a\over 2}.
\label{eq:RodweinormanJ}
\eeq
Similarly, we could leave $\bfOmega_b$ in the product expansion (\ref{eq:expproduct}) at a finite value
and consider infinitesimal left variations of $\bfOmega_a\rightarrow {\bf 0}$. This naturally leads to the
complementary left Rodrigues Jacobian
\beq
\mathsf{J}_{\omega_l} (\bfOmega_b) = (1-y)\, \mathsf{I}+ y\, {\bf r}_b{\bf r}_b^T - {1\over 2}[\bfOmega_b]_\times,
\; 1-y = {\Theta_b\over 2}\tan^{-1}{\Theta_b\over 2}.
\eeq
It shall be emphasised that the Jacobians above don't have singularities and are invertible, since all quantities
are well defined for $\Theta\leq\pi$ and $\mathrm{det}( \mathsf{J}_{\omega_r}) = (1-x)^2 + \Theta_a^2/4$.

%% -----------------------------------------------------------------------
\subsection{$SE(3)$ Composition Jacobians}
\label{app:SE3composition}
%% -----------------------------------------------------------------------
In close analogy to the previous sections, a composition in $SE(3)$ can be analyzed
in the {\em additive} framework,
\beq
(\mathsf{R}(\bfOmega_c),{\bf T}_c)=(\mathsf{R}_a,{\bf T}_a)\cdot
(\mathsf{R}_b,{\bf T}_b)\, .
\eeq

\subsubsection{Composing finite motions}
With the group composition (\ref{eq:groupcomp}), we can formally write
\beq
\ln \, (\mathsf{R}(\bfOmega_c),{\bf T}_c) = [\ln \mathsf{R}_a\mathsf{R}_b,\;
\mathsf{S}^{-1}(\bfOmega_c) (\mathsf{R}_a{\bf T}_b + {\bf T}_a) \, ].
\label{eq:addfiniteSE3}
\eeq

\subsubsection{Composing a finite with an infinitesimal motion}
The addition formula (\ref{eq:addfiniteSE3}) can be analyzed in the limit $t\rightarrow 0$ where the right composition is $\ln (\mathsf{R}_b,{\bf T}_b) = [t \bfomega_b, t {\bf v}_b]$. It is easy to show that the following must hold to order ${\cal O}(t)$,
\begin{multline}
    [\bfOmega_c,{\bf t}_c] = [\bfOmega_a,{\bf t}_a] + \\
    [\mathsf{J}_{\omega_r}(\bfOmega_a)\bfomega_b,\,
    \mathsf{S}^{-1}(\bfOmega_a)\mathsf{R}(\bfOmega_a){\bf v}_b + \mathsf{J}_{t_r}(\bfOmega_a,{\bf t}_a)\bfomega_b].
    \label{eq:SE3vectoraddition}
\end{multline}

\subsubsection{The $\mathsf{S}^{-1}$ identity and coupling Jacobian}
In (\ref{eq:rodriguesapp}) we only gave the expression for $\mathsf{S}$. We can use the
properties of the finite $\mathsf{H}$ algebra with its characteristic polynomial $\mathsf{H}(\mathsf{H}^2+\mathsf{I})=0$
to obtain an expression for $\mathsf{S}^{-1}(\bfOmega_a)$. As an interesting fact, we meet the Rodrigues Jacobians again
\begin{multline}
    \mathsf{S}^{-1}(\bfOmega_a) = (1-x)\mathsf{I} + x\, {\bf r}_a{\bf r}_a^T - {1\over 2}[\bfOmega_a]_\times \\
    = \mathsf{J}_{\omega_r}(-\bfOmega_a) = \mathsf{J}_{\omega_l}(\bfOmega_a)\;.
    \label{eq:Sinvidentity}
\end{multline}
Further, one can show that
\beq
\mathsf{S}^{-1}(\bfOmega_a)\mathsf{R}(\bfOmega_a) = \mathsf{R}(\bfOmega_a)\mathsf{S}^{-1}(\bfOmega_a) =
\mathsf{J}_{\omega_r}(\bfOmega_a),
\label{eq:SinvRidentity}
\eeq
A more detailed calculation for the coupling Jacobian $\mathsf{J}_{t_r}$ yields
\beq
\mathsf{J}_{t}(\bfOmega_a,{\bf t}_a) = {\partial [\mathsf{S}^{-1}(\bfOmega_a){\bf T}_a] \over
\partial \bfOmega_a} \, \mathsf{J}_{\omega_r}(\bfOmega_a),\quad {\bf T}_a = \mathsf{S}(\bfOmega_a) \, {\bf t}_a.
\label{eq:couplingJacobian}
\eeq
We can show that
\beqar
{\partial [\mathsf{S}^{-1}(\bfOmega_a){\bf T}_a] \over
    \partial \bfOmega_a} & = &
{1\over 2}[{\bf t}_a]_\times+ {x\over\Theta_a}
\left( {\bf r}_a{\bf t}_a^T- t_a \;\mathsf{H}^2_a + {\bf t}_a{\bf r}_a^T \right) + \nonumber \\
&& {t_a\over\Theta_a} (\mathsf{S}(\bfOmega_a)-\mathsf{S}^{-1}(-\bfOmega_a)),
\label{eq:dSinvTdOmega}
\eeqar
with $ t_a = {\bf r}_a\cdot{\bf t}_a$, which finishes the discourse on the $SE(3)$ composition Jacobians.

\subsection{Additive versus compositive derivative}
\label{app:addvscomp}
The Rodrigues formula (\ref{eq:rodriguesapp}) allows differentiation with respect to $\bfOmega$ directly,
\beq
\mathsf{R}(\bfOmega_0+\Delta\bfomega) \approx
\mathsf{R}(\bfOmega_0) + \left. {\partial\mathsf{R}\over \partial\Omega_\mu} \right|_{\bfOmega_0}
\Delta\omega_\mu = \mathsf{R}(\bfOmega_0) ( \mathsf{I} + [\delta\bfOmega]_\times ).
\label{eq:rodriguesder}
\eeq
We can call $\partial\mathsf{R}/ \partial\Omega_\mu$ the additive derivative, whereas
$[\delta\bfOmega]_\times = \left.{d\over dt}\right|_{t=0} \exp( t [\delta\bfOmega]_\times )$
can be termed the compositive ({\em incremental}) derivative.

In accordance with the vector addition over $SO(3)$ (\ref{eq:addSO3finite}), we can show that
\beq
\delta\bfOmega \; = \; \mathsf{J}^{-1}_{\omega_r}(\bfOmega_0)\, \Delta\bfomega
\label{eq:SO3localglobal}
\eeq
or explicitly
\beq
\delta{\bfOmega}=
{\sin\Theta_0\over\Theta_0} \Delta\bfomega + {1-\cos\Theta_0\over\Theta_0}\; \Delta\bfomega\times{\bf r}_0 +
k\,(1-{\sin\Theta_0\over\Theta_0}) \; {\bf r}_0
\nonumber
\eeq
with $k={\bf r}_0\cdot\Delta\bfomega$ and the usual definition $\Theta_0=|\bfOmega_0|$,
${\bf r}_0 = \bfOmega_0/\Theta_0$.

%% ====================================================================
\section{Second-Order Propagation Formulas}
\label{app:secondorder}
%% ====================================================================

We derive the second-order corrections to mean and covariance
for the composition $h = g_0 \cdot \Delta g \cdot f_0 \cdot \Delta f$,
where $\Delta g = \exp([\bfomega_g, {\bf v}_g])$ and
$\Delta f = \exp([\bfomega_f, {\bf v}_f])$ are zero-mean
perturbations with covariances $\sfSigma{gg}$ and $\sfSigma{ff}$ respectively.

\subsection{Mean correction from BCH curvature}

Using the finite composition formula~(\ref{eq:addSO3finite}) rather than
the truncated BCH series, the composed rotation in exponential coordinates
can be expanded to second order as
\beq
\bfOmega_c = \bfOmega_a + \bfOmega_b + \frac{1}{2}\bfOmega_a \times \bfOmega_b
+ \frac{1}{12}\left(\bfOmega_a \times (\bfOmega_a \times \bfOmega_b)
+ \bfOmega_b \times (\bfOmega_b \times \bfOmega_a)\right) + \cdots
\label{eq:BCHexpanded}
\eeq
Taking the expectation with $\langle \bfOmega_a \rangle = 0$,
$\langle \bfOmega_b \rangle = 0$:
\beq
\langle \bfOmega_c \rangle = \frac{1}{2}\langle \bfOmega_a \times \bfOmega_b \rangle
= \frac{1}{2} \sum_{\alpha\beta} \varepsilon_{\cdot\alpha\beta}\,
\langle \omega_{a,\alpha}\, \omega_{b,\beta} \rangle.
\label{eq:meanshift_rotation}
\eeq
For the translational component in $SE(3)$, the coupling through $\mathsf{S}^{-1}$
introduces additional terms. Using~(\ref{eq:SE3vectoraddition}) expanded to second order,
the full $SE(3)$ mean correction is
\beq
\delta\mu^{(2)} = \frac{1}{2}
\begin{pmatrix}
    \varepsilon_{i\alpha\beta}\, [\sfSigma{gf}]_{\alpha\beta}^{\omega\omega} \\[4pt]
    \varepsilon_{i\alpha\beta}\, [\sfSigma{gf}]_{\alpha\beta}^{\omega v}
    + [\mathsf{Q}(\bfOmega_0)]_{ij}\, [\sfSigma{ff}]_{j}^{\omega\omega}
\end{pmatrix},
\label{eq:fullmeancorrection}
\eeq
where $\mathsf{Q}(\bfOmega_0)$ encodes the Hessian of $\mathsf{S}^{-1}(\bfOmega)\mathsf{R}(\bfOmega)$
evaluated at the base point, and the superscripts denote the rotational ($\omega$)
and translational ($v$) blocks of the cross-covariance.

\subsection{Covariance correction}

The second-order covariance correction arises from the
quadratic term $Q_i = \frac{1}{2}\varepsilon_{i\alpha\beta}\,
    \omega_{a,\alpha}\,\omega_{b,\beta}$ in the BCH expansion.
Writing $\bfOmega_c = L + Q + \cdots$ where
$L_i = \omega_{a,i} + \omega_{b,i}$ is the linear part,
the covariance picks up a correction from
$\langle Q_i Q_j \rangle - \langle Q_i \rangle \langle Q_j \rangle$.

\subsubsection{Rotational block ($3\times 3$)}
Since $Q_i$ mixes components from two distinct random
vectors $\bfomega_a$ (from $\Delta g$) and $\bfomega_b$
(from $\Delta f$), the indices must be tracked separately:
\beq
\Delta\Sigma^{(2),\omega\omega}_{ij} = \frac{1}{4}
\sum_{\alpha\beta\gamma\delta}
\varepsilon_{i\alpha\beta}\,\varepsilon_{j\gamma\delta}\,
\mathcal{C}_{\alpha\beta\gamma\delta},
\label{eq:covcorrection_rotrot}
\eeq
where the fourth-moment tensor is
\beq
\mathcal{C}_{\alpha\beta\gamma\delta} =
\langle \omega_{a,\alpha}\,\omega_{b,\beta}\,
\omega_{a,\gamma}\,\omega_{b,\delta} \rangle
- \langle \omega_{a,\alpha}\,\omega_{b,\beta} \rangle
\langle \omega_{a,\gamma}\,\omega_{b,\delta} \rangle.
\eeq
For Gaussian perturbations, the fourth-order cumulant
vanishes identically. The correction comes entirely from
the factorization of the fourth moment via Isserlis' theorem:
\begin{multline}
    \langle \omega_{a,\alpha}\,\omega_{b,\beta}\,
    \omega_{a,\gamma}\,\omega_{b,\delta} \rangle =
    [\sfSigma{aa}]_{\alpha\gamma}\,[\sfSigma{bb}]_{\beta\delta} \\
    + [\sfSigma{ab}]_{\alpha\delta}\,[\sfSigma{ba}]_{\beta\gamma}
    + [\sfSigma{ab}]_{\alpha\beta}\,[\sfSigma{ab}]_{\gamma\delta}.
    \label{eq:isserlis}
\end{multline}
The last term cancels with the subtracted product, giving
the closed-form correction:
\beq
\boxed{
    \Delta\Sigma^{(2),\omega\omega}_{ij} = \frac{1}{4}
    \sum_{\alpha\beta\gamma\delta}
    \varepsilon_{i\alpha\beta}\,\varepsilon_{j\gamma\delta}
    \Big([\sfSigma{aa}]_{\alpha\gamma}\,[\sfSigma{bb}]_{\beta\delta}
    + [\sfSigma{ab}]_{\alpha\delta}\,[\sfSigma{ba}]_{\beta\gamma}\Big)
}
\label{eq:covcorrection_explicit}
\eeq
When the perturbations are independent ($\sfSigma{ab} = 0$),
this simplifies to
\beq
\Delta\Sigma^{(2),\omega\omega}_{ij} = \frac{1}{4}
\varepsilon_{i\alpha\beta}\,\varepsilon_{j\gamma\delta}\,
[\sfSigma{aa}]_{\alpha\gamma}\,[\sfSigma{bb}]_{\beta\delta},
\eeq
which can be evaluated using the identity
$\varepsilon_{i\alpha\beta}\varepsilon_{j\gamma\delta}
    A_{\alpha\gamma}B_{\beta\delta} =
    (\mathrm{tr}\,\mathsf{A})(\mathrm{tr}\,\mathsf{B})\delta_{ij}
    - (\mathrm{tr}\,\mathsf{A})B_{ij} - A_{ij}(\mathrm{tr}\,\mathsf{B})
    + A_{ik}B_{kj} + A_{jk}B_{ki}$.

\subsubsection{Rotation-translation cross block ($3\times 3$)}
The translational component $t_{c,i}$ receives a second-order
contribution through the coupling
$\mathsf{S}^{-1}(\bfOmega)\mathsf{R}(\bfOmega)$,
expanded around the base point. Using the
SE(3) vector addition~(\ref{eq:SE3vectoraddition}),
the cross-covariance correction is
\begin{multline}
    \Delta\Sigma^{(2),\omega t}_{ij} = \frac{1}{2}
    \varepsilon_{i\alpha\beta}\Big(
    [\sfSigma{aa}]_{\alpha\gamma}\,
    [\mathsf{J}_{\omega_r}]_{\gamma k}\,
    [\sfSigma{bb}^{vv}]_{\beta j} \\
    + [\sfSigma{ab}^{\omega v}]_{\alpha j}\,
    [\sfSigma{ba}^{\omega\omega}]_{\beta\gamma}\,
    [\mathsf{J}_t]_{\gamma k}\Big)
    + \mathcal{O}(\sigma^5),
    \label{eq:covcorrection_rottrans}
\end{multline}
where $[\sfSigma{}^{vv}]$, $[\sfSigma{}^{\omega v}]$
denote the translational and cross blocks of the
$6\times 6$ covariance, and $\mathsf{J}_t$ is the
coupling Jacobian from~(\ref{eq:couplingJacobian}).

\subsubsection{Translational block ($3\times 3$)}
The pure translational correction receives contributions
from: (i)~the $\mathsf{J}_t$ coupling acting on the
rotational perturbation, and (ii)~the Hessian of
$\mathsf{S}^{-1}\mathsf{R}$ at the base point.
To leading order in $\sigma^4$:
\begin{multline}
    \Delta\Sigma^{(2),tt}_{ij} =
    [\mathsf{Q}]_{ik}\,[\sfSigma{aa}^{\omega\omega}]_{kl}\,
    [\mathsf{Q}]^T_{lj} \\
    + \frac{1}{4}\varepsilon_{i'\alpha\beta}\,
    \varepsilon_{j'\gamma\delta}\,
    [\mathsf{J}_t]_{ii'}\,[\mathsf{J}_t]_{jj'}\,
    \mathcal{C}_{\alpha\beta\gamma\delta}
    + \mathcal{O}(\sigma^5),
    \label{eq:covcorrection_transtrans}
\end{multline}
where $\mathsf{Q}(\bfOmega_0)$ is the Hessian matrix from
the mean correction~(\ref{eq:fullmeancorrection}) and
$\mathcal{C}_{\alpha\beta\gamma\delta}$ is the same
Isserlis-reduced tensor as in~(\ref{eq:covcorrection_explicit}).

\begin{remark}
    The key insight is that for Gaussian perturbations, no
    higher-order moments are needed---the entire correction
    is determined by the second-order covariance blocks
    $\sfSigma{aa}$, $\sfSigma{bb}$, $\sfSigma{ab}$,
    combined with the structure constants $\varepsilon_{ijk}$
    and the coupling Jacobian $\mathsf{J}_t$. The correction
    vanishes on commutative groups (where $\varepsilon_{ijk}=0$)
    and on $SO(3)$ only the rotational
    block~(\ref{eq:covcorrection_explicit}) survives.
\end{remark}

%% ====================================================================
\section{Saddlepoint Approximation Details}
\label{app:saddlepoint}
%% ====================================================================

\subsection{Moment generating function for projective observations}

For a single landmark ${\bf x}$ observed from camera $f$ with projective
model $\pi(f\star{\bf x}) = [u, v, 1]$ where $u = x_1'/x_3'$, $v = x_2'/x_3'$,
the log-likelihood as a function of ${\bf x}$ (with $f$ fixed) is
\beq
\ell({\bf x}) = -\frac{1}{2}({\bf z} - \pi(f\star{\bf x}))^T \sfSigma{zz}^{-1}
({\bf z} - \pi(f\star{\bf x})).
\eeq
The moment generating function of the integrand in the marginalization~(\ref{eq:marginalize}) is
\beq
K({\bf s}) = \ell({\bf x}_{\rm opt}) + \frac{1}{2}{\bf s}^T \mathsf{H}^{-1} {\bf s}
+ \frac{1}{6}\sum_{abc} \kappa_{abc}\, s_a s_b s_c + \cdots
\label{eq:cumulantexpansion}
\eeq
where $\mathsf{H} = -\nabla^2 \ell |_{{\bf x}_{\rm opt}}$ is the Hessian
(positive definite at the optimum) and $\kappa_{abc}$ are the third cumulants.

\subsection{Third cumulants of the projective model}

The third derivatives of $\pi$ at ${\bf x}' = f\star{\bf x}$ are:
\beq
\frac{\partial^3 u}{\partial x_a' \partial x_b' \partial x_c'} =
\frac{\partial^3}{\partial x_a' \partial x_b' \partial x_c'}\left(\frac{x_1'}{x_3'}\right),
\eeq
which are rational functions with poles at $x_3' = 0$ (behind the camera).
The non-vanishing components are:
\beqar
\frac{\partial^3 u}{\partial x_1' \partial x_3'^2} &=& \frac{2}{x_3'^3}, \quad
\frac{\partial^3 u}{\partial x_3'^3} = \frac{6\, u}{x_3'^3}, \nonumber \\
\frac{\partial^3 v}{\partial x_2' \partial x_3'^2} &=& \frac{2}{x_3'^3}, \quad
\frac{\partial^3 v}{\partial x_3'^3} = \frac{6\, v}{x_3'^3}.
\label{eq:thirdderivs}
\eeqar

The third cumulants in the original landmark coordinates ${\bf x}$
are obtained by chain rule through $\mathsf{R}_f$:
\beq
\kappa_{abc} = \sum_{a'b'c'} \frac{\partial^3 \ell}{\partial x_{a'}' \partial x_{b'}' \partial x_{c'}'}
\, R_{a'a}\, R_{b'b}\, R_{c'c}.
\label{eq:thirdcumulants}
\eeq

\subsection{Saddlepoint correction to Laplace}

The saddlepoint approximation to the marginal contribution of
landmark $j$ is obtained by solving the saddlepoint equation
$K'({\bf s}_0) = 0$ and evaluating:
\beq
\int e^{\ell({\bf x})} d{\bf x} \approx
e^{\ell({\bf x}_{\rm opt})} \, (2\pi)^{3/2} |\mathsf{H}|^{-1/2}
\left(1 + \frac{1}{8}\sum_{abcd} \kappa_{abcd}\, [\mathsf{H}^{-1}]_{ab}[\mathsf{H}^{-1}]_{cd}
+ \frac{5}{24}\sum_{abc} \kappa_{abc}^2 \, [\mathsf{H}^{-3}]_{abc}\right),
\label{eq:saddlepointfull}
\eeq
where $\kappa_{abcd}$ are the fourth cumulants and
$[\mathsf{H}^{-3}]_{abc} = \sum_{def} [\mathsf{H}^{-1}]_{ad}[\mathsf{H}^{-1}]_{be}[\mathsf{H}^{-1}]_{cf}$.

The Laplace approximation retains only the leading term
(the first factor in~(\ref{eq:saddlepointfull})). The correction
in parentheses is $\mathcal{O}(1/n)$ for $n$ observations
and becomes significant when:
(a)~the depth $x_3'$ is comparable to the baseline,
(b)~few cameras observe the landmark, or
(c)~the landmark is near the image boundary where
$u$ or $v$ are large.

The corrected marginal log-posterior over camera poses is then
\beq
\ln p(f | {\bf z}) = \sum_j \left[\ell_j({\bf x}^j_{\rm opt}(f))
- \frac{1}{2}\ln|\mathsf{H}_j(f)| + \ln(1 + \delta_j^{\rm SP}(f))\right]
+ \ln p(f),
\label{eq:correctedmarginal}
\eeq
where $\delta_j^{\rm SP}$ is the saddlepoint correction term from~(\ref{eq:saddlepointfull}).
This expression is optimized over $f \in SE(3)$ using the Gauss-Newton iteration
with compositive updates $f_{t+1} = f_t \cdot \exp(\Delta{\bf f}_t)$ as described
in Section~\ref{sec:saddlepoint}.

%% ====================================================================
%% BIBLIOGRAPHY
%% ====================================================================
\begin{thebibliography}{99}

    \bibitem{pennec97}
    X.~Pennec and J.-P.~Thirion,
    ``A framework for uncertainty and validation of 3-D registration
    methods based on points and frames,''
    \emph{Int. J. Computer Vision}, vol.~25, no.~3, pp.~203--229, 1997.

    \bibitem{barrau2017invariant}
    A.~Barrau and S.~Bonnabel,
    ``The invariant extended Kalman filter as a stable observer,''
    \emph{IEEE Trans. Autom. Control}, vol.~62, no.~4, pp.~1797--1812, 2017.

    \bibitem{sola2018micro}
    J.~Sol\`{a}, J.~Deray, and D.~Atchuthan,
    ``A micro Lie theory for state estimation in robotics,''
    \emph{arXiv:1812.01537}, 2018.

    \bibitem{barfoot2024}
    T.~D.~Barfoot,
    \emph{State Estimation for Robotics}, 2nd ed.
    Cambridge University Press, 2024.

    \bibitem{forster2017manifold}
    C.~Forster, L.~Carlone, F.~Dellaert, and D.~Scaramuzza,
    ``On-manifold preintegration for real-time visual-inertial odometry,''
    \emph{IEEE Trans. Robotics}, vol.~33, no.~1, pp.~1--21, 2017.

    \bibitem{chirikjian2012stochastic}
    G.~S.~Chirikjian,
    \emph{Stochastic Models, Information Theory, and Lie Groups,
        Volume~2: Analytic Methods and Modern Applications}.
    Birkh\"{a}user, 2012.

    \bibitem{ye2024uncertainty}
    J.~Ye and G.~S.~Chirikjian,
    ``Uncertainty propagation and Bayesian fusion on unimodular
    Lie groups from a parametric perspective,''
    \emph{arXiv:2401.03425}, 2024.

    \bibitem{wang2008nonparametric}
    Y.~Wang and G.~S.~Chirikjian,
    ``Nonparametric second-order theory of error propagation
    on motion groups,''
    \emph{Int. J. Robotics Research}, vol.~27, no.~11--12,
    pp.~1258--1273, 2008.

    \bibitem{bourmaud2013discrete}
    G.~Bourmaud, R.~M\'{e}gret, A.~Giremus, and Y.~Berthoumieu,
    ``Discrete extended Kalman filter on Lie groups,''
    \emph{Proc. EUSIPCO}, pp.~1--5, 2013.

    \bibitem{triggs-bundle}
    B.~Triggs, P.~McLauchlan, R.~Hartley, and A.~Fitzgibbon,
    ``Bundle adjustment---a modern synthesis,''
    in \emph{Vision Algorithms: Theory and Practice},
    Springer, 2000, pp.~298--375.

    \bibitem{Horn87}
    B.~K.~P.~Horn,
    ``Closed-form solution of absolute orientation using unit
    quaternions,''
    \emph{J. Opt. Soc. Am. A}, vol.~4, no.~4, pp.~629--642, 1987.

    \bibitem{Wei-Norman}
    J.~Wei and E.~Norman,
    ``On the global representations of the solutions of linear
    differential equations as a product of exponentials,''
    \emph{Proc. Amer. Math. Soc.}, vol.~15, pp.~327--334, 1964.

    \bibitem{kuehnel2004}
    F.~O.~Kuehnel,
    ``Bayesian estimation of nonlinear parameters on $SE(3)$
    Lie group,''
    in \emph{Bayesian Inference and Maximum Entropy Methods
        in Science and Engineering} (MaxEnt 2004),
    R.~Fischer, R.~Preuss, and U.~von Toussaint, Eds.,
    AIP Conf.\ Proc., vol.~735, pp.~176--186, 2004.
    DOI:~10.1063/1.1835212.

    \bibitem{kuehnel2005}
    F.~O.~Kuehnel,
    ``Local frame junction trees in SLAM,''
    \emph{AIP Conf.\ Proc.}, vol.~803, pp.~318--329, 2005.
    DOI:~10.1063/1.2149810.

    \bibitem{kuehnel2008}
    F.~O.~Kuehnel,
    ``Robust Bayesian estimation of nonlinear parameters on
    $SE(3)$ Lie group from noisy measurements,''
    Tech.\ Rep., NASA Ames / USRA-RIACS, 2008.

\end{thebibliography}

\end{document}
